\section{Recap}
\begin{frame}{}
    \LARGE Continuous Control: \textbf{Recap (Day 4)}
\end{frame}

\begin{frame}{Recap: What Day 4 Established}
\begin{itemize}
    \item \textbf{Actor-Critic}: an actor $\pi_\theta(a|s)$ updated by the policy gradient, plus a critic $V_\phi(s)$ trained by TD learning to supply a low-variance advantage estimate $\hat{A}_t$:
    $$\nabla_\theta \mathcal{J}(\theta) = \mathbb{E}_t\!\left[ \nabla_\theta \log \pi_\theta(a_t|s_t)\, \hat{A}_t \right]$$
    \item The family tree ran: A2C/A3C $\rightarrow$ +GAE $\rightarrow$ TRPO $\rightarrow$ \textbf{PPO} $\rightarrow$ \textbf{GRPO} (critic-free, group-mean baseline --- the current LLM-RL default)
    \item \textbf{Key limitation (the one we attack today):} every method in that tree is \emph{on-policy} --- it \emph{discards all experience after each update}. Reusing a replay buffer is the central win today
    \item \emph{Secondary:} the actor is always \emph{stochastic}. For a physical actuator you eventually just want one number per joint --- no distribution to sample or integrate over
\end{itemize}
\end{frame}

\begin{frame}{Recap: What Today Delivers}
\begin{itemize}
    \item \textbf{Today we go off-policy and continuous.}
    \item \textbf{DPG} \rref{1}: a policy-gradient update rule for a \emph{deterministic} actor $\mu_\theta(s)$ over continuous action spaces
    \item \textbf{DDPG} \rref{2}: DPG made deep, by importing the DQN \rref{3} tricks you already know --- replay buffer and target networks
    \item The critic becomes an \emph{action}-value $Q_\phi(s,a)$, and the actor's gradient flows \emph{through the critic}: $\nabla_\theta Q_\phi(s, \mu_\theta(s))$
    \item We close with \textbf{TD3} \rref{4}, the direct fix for DDPG's overestimation and brittleness --- and the bridge to Day 6 (SAC)
    \vspace{0.4em}
    \item[] \footnotesize \textcolor{red}{Scope note:} PPO already handles continuous actions \emph{on-policy}. ``Continuous Control I/II'' is specifically the \emph{off-policy} branch: deterministic policies (DDPG/TD3, today) and max-entropy stochastic policies (SAC, Day 6).
\end{itemize}
\end{frame}

\begin{frame}{Notation: What Carries Over, What Is New}
\vspace{-0.2em}
\begin{center}
\small
\renewcommand{\arraystretch}{1.25}
\begin{tabular}{lll}
\hline
\textbf{Symbol} & \textbf{Meaning} & \textbf{Where you saw it} \\
\hline
$\mathcal{J}(\theta)$        & expected discounted return (the objective)   & Days 3--4 \\
$\theta$                     & \textbf{actor} parameters                    & Days 3--4 \\
$\phi$                       & \textbf{critic} parameters                   & Day 4 \\
$\theta^{-},\ \phi^{-}$      & \emph{target} network parameters             & Day 2 (DQN) \\
$\mathcal{D}$                & replay buffer                                & Day 2 (DQN) \\
$\pi_\theta(a|s)$            & \emph{stochastic} actor                      & Days 3--4 \\
$\mu_\theta(s)$              & \emph{deterministic} actor \; \textcolor{red}{(new today)} & Day 3, as the Gaussian mean \\
$Q_\phi(s,a)$                & \emph{action}-value critic \; \textcolor{red}{(new today)} & Day 4's critic was $V_\phi(s)$ \\
\hline
\end{tabular}
\end{center}
\vspace{0.5em}
{\footnotesize
\textbf{Why $\mu$ for the actor?} Day 3 already wrote the continuous-action policy as a Gaussian,
$\pi_\theta(a|s) = \mathcal{N}\!\big(\mu_\theta(s),\, \sigma_\theta(s)^2\big)$.
Today we keep the mean and throw the noise away: the actor \emph{is} $\mu_\theta(s)$. Same letter, same network --- we have simply deleted $\sigma_\theta$. Keeping $\pi$ and $\mu$ distinct also matters on Day 6, where SAC puts a stochastic actor back.}
\end{frame}
